% Propietario conservado: /Users/ruben/Documents/ChatGPT/jueces y controles/REVISION_ARGUMENTAL_APERTURAS_CIERRES_20260915/II/manuscript_es/sections/definicion_gravedad_rev06.tex
\subsection{Définition structurelle et fonction de la constante gravitationnelle}
\label{ii:sec:ii-definicion-gravedad}

\begin{definition}[Section gravitationnelle contragrédiente de l'action]
Pour une section positive \(\hbar_a\), une vitesse
\(c_{\rm int}>0\) et le rayon chronologique \(R_{108}\) de la
section~\ref{ii:sec:gravity}, la section gravitationnelle circulaire est
\begin{equation}
 G_a=\frac{c_{\rm int}^{3}R_{108}^{2}}{\hbar_a},
 \qquad R_{108}=\frac{54}{\pi}\ell_0,
 \qquad \ell_0=c_{\rm int}t_0.
 \label{ii:eq:ii-definicion-gravedad-seccion}
\end{equation}
C'est la section réciproque de l'action qui réalise la coïncidence
entre le rayon gravitationnel réduit du cycle et son rayon
chronologique, dans la convention \(r_g=Gm/c_{\rm int}^2\).
\end{definition}

La proposition de coïncidence des rayons de la
section~\ref{ii:sec:gravity} démontre son unicité dans la famille
\(G_a(\vartheta)=\vartheta c_{\rm int}^{5}t_0^2/\hbar_a\).
La condition circulaire fixe \(\vartheta=(54/\pi)^2\) et conduit
à \eqref{ii:eq:ii-definicion-gravedad-seccion}. Cette condition et les
sections dimensionnelles positives font partie de la réalisation.
En les maintenant communes aux deux orientations, on obtient
\begin{equation}
 \frac{G_{\mathrm{pre}}}{G_{\mathrm{ret}}}
 =R_{\mathrm{act}}^{-1},
 \qquad
 \frac{G_a\hbar_a}{c_{\rm int}^{3}}
 =R_{108}^{2}=\ell_P^2.
 \label{ii:eq:ii-definicion-gravedad-area}
\end{equation}
La preuve est la substitution directe de
\eqref{ii:eq:ii-definicion-gravedad-seccion} ; le changement d'action
est compensé par le changement réciproque de \(G_a\).

Cette structure attribue à la gravitation une fonction précise :
transporter la section d'action vers une unité géométrique commune.
Dans \eqref{ii:eq:holst}, \(G\) fixe la normalisation du couplage ;
dans l'opérateur d'aire, \(G\hbar/c^3\) fixe l'unité aréale.
La réciprocité conserve cette unité tandis que change la lecture
orientée de l'action. L'état réduit et ses occupations apportent
ensuite l'information et l'intrication associées à la frontière.

Le contenu tensoriel du couplage est conservé dans
\(\mathbb G_{5,a}=G_aP^{\mathrm{ico}}_5\), construit dans
\ref{ii:sec:ii-pentadica}. La section \(G_a\) détermine sa grandeur commune ;
le projecteur sélectionne le porteur à cinq composantes et \(\Gamma_5\)
transporte son orientation vers l'espace des tenseurs symétriques sans trace.
La base de \eqref{ii:ii:pent:eq:gravity-helicity2-real}--\eqref{ii:ii:pent:eq:gravity-helicity0-real}
explicite les cinq composantes. La projection \(\Lambda(n)\) définit
ensuite la réduction transversale. Grandeur, porteur et réduction sont des
opérations liées, dont les domaines sont différents au sein de cette réalisation.
